\clearpage
\begin{summarybox}
  \textbf{\textcolor{CSummary}{一分钟记住}}
  \[
     \boxed{\ \text{横坐标}\times k\ (k>0)
     \quad\Longrightarrow\quad y=f\!\left(\frac{x}{k}\right)\ }
  \]
  换成 $y=f(\omega x)$ 的说法，就是\textbf{横坐标乘 $1/\omega$}（$\omega>0$）。

  \medskip
  \textbf{\textcolor{CSummary}{一句话理解}}
  \textbf{同样的高度，要用新的横坐标找到原来的输入。} 例如 $f(2x)$ 中，新横坐标 $1$ 对应原来输入 $2$，所以图像横向压成一半。

  \medskip
  \textbf{\textcolor{CSummary}{三个例子}}
  \begin{itemize}[leftmargin=*,itemsep=0.28em,topsep=0.3em]
    \item 横坐标压成 $1/3$：$f(x)\to f(3x)$。
    \item 横坐标拉成 $2$ 倍：$f(x)\to f(x/2)$。
    \item 横坐标不变：$f(x)\to f(x)$。
  \end{itemize}

  \medskip
  \textbf{\textcolor{CSummary}{考试最快做法}}
  看清题目说的是\textbf{横坐标变化}还是\textbf{括号内系数}；写新式子后，任选一个点，检查\textbf{高度是否保持不变}。遇到平移和伸缩混合，严格按题目顺序处理。

  \medskip
  \textbf{\textcolor{CSummary}{记忆口诀}}
  \textbf{横向变化看倒数，里面乘大图变窄，里面乘小图变宽。}
\end{summarybox}
